% figure-1-examples-of-cpos.tex
%
% Source:   Carl A. Gunter and Dana S. Scott, "Semantic Domains" (1990).
%           ScottDomains/papers/Gunter Scott 1990.pdf
% Physical page: 5           Printed page: 4
% Section:  2.1 "Cpo's and the Fixed Point Theorem" (the figure sits at the
%           top of printed page 4; the prose introducing it is on printed
%           page 3 / physical page 4).
% Caption (verbatim):  "Figure 1: Examples of cpo's."
%
% Introducing sentences (printed page 3 / physical page 4):
%   "The truth value cpo T is the poset which has three distinct points,
%    \bot, true, false, where \bot \sqsubseteq true and \bot \sqsubseteq false
%    (see Figure 1)."
%   "To get a cpo, we need to add a \"bottom\" element to N. The result is a
%    cpo N_\bot which is pictured in Figure 1."
%   "To get a cpo, one needs to add a top element to get the cpo \omega^\top
%    pictured in Figure 1."
%
% Reading of the three parts:
%   T          the truth value cpo: \bot below the two incomparable points
%              true and false.
%   N_\bot     the natural numbers with the discrete ordering, with a bottom
%              element adjoined; the row of maximal points continues to the
%              right, so one edge is drawn truncated and followed by \cdots.
%   \omega^\top  the ordinal \omega, an ascending chain with no least upper
%              bound, with a top element adjoined above the whole chain; the
%              chain's continuation is a truncated vertical followed by \vdots.
%
% Every point in this figure is drawn hollow in the paper - there is no
% distinguished pair here, unlike Figure 3 - so every vertex uses the "open"
% style.
%
% Geometry was measured from 600 dpi crops of physical page 5 and mapped to
% centimetres by x = (page pixel - 715)/200 and y = (2503 - page pixel)/200,
% so the three parts keep the paper's relative placement and scale, including
% the omega^\top column sitting slightly higher than the other two.
%
% Style contract (r0039) with the two orchestrator-settled corrections:
% the contract's "\usetikzpicture{}" line is a typo and is dropped, and
% "solid" is 5pt against "open" at 4pt (no solid vertex occurs in this
% figure). The option block also carries "obj", "arrow" and "darrow",
% agent2's three keys for commutative diagrams, so every file in this round's
% set has an identical preamble; this figure is a Hasse diagram and uses none
% of the three.

\documentclass[border=6pt]{standalone}
\usepackage{tikz}
\begin{document}
\begin{tikzpicture}[
  x=1cm, y=1cm,
  every node/.style={inner sep=1pt},
  open/.style={circle, draw, fill=white, minimum size=4pt, inner sep=0pt},
  solid/.style={circle, draw, fill=black, minimum size=5pt, inner sep=0pt},
  edge/.style={draw, line width=0.4pt},
  obj/.style={inner sep=3pt},
  arrow/.style={draw, line width=0.4pt, ->},
  darrow/.style={draw, line width=0.4pt, ->, dash pattern=on 4pt off 3pt}]

  %% ---------------- T, the truth value cpo ----------------
  \node[open] (tl) at (0.000, 2.23) {};
  \node[open] (tr) at (4.485, 2.23) {};
  \node[open] (tb) at (2.215, 0.00) {};
  \draw[edge] (tb) -- (tl);
  \draw[edge] (tb) -- (tr);
  \node at (2.20, -1.35) {$\mathsf{T}$};

  %% ---------------- N_\bot, the flat natural numbers ----------------
  \node[open] (n1) at (7.475, 2.23) {};
  \node[open] (n2) at (8.985, 2.23) {};
  \node[open] (n3) at (10.495, 2.23) {};
  \node[open] (nb) at (9.745, -0.03) {};
  \draw[edge] (nb) -- (n1);
  \draw[edge] (nb) -- (n2);
  \draw[edge] (nb) -- (n3);
  % the row of maximal points continues to the right
  \draw[edge] (nb) -- (11.825, 2.085);
  \node at (12.235, 2.23) {$\cdots$};
  \node at (9.745, -1.35) {$\mathsf{N}_{\bot}$};

  %% ---------------- \omega^\top, the ordinal \omega with a top ----------------
  \node[open] (w0) at (16.85, 0.250) {};
  \node[open] (w1) at (16.85, 1.615) {};
  \node[open] (w2) at (16.85, 2.975) {};
  \draw[edge] (w0) -- (w1);
  \draw[edge] (w1) -- (w2);
  % the chain continues upward, then the adjoined top element
  \draw[edge] (w2) -- (16.85, 4.315);
  \node at (16.85, 5.055) {$\vdots$};
  \node[open] (wt) at (16.85, 5.975) {};
  \node at (16.85, -1.35) {$\omega^{\top}$};

\end{tikzpicture}
\end{document}
